\documentclass[11pt,a4paper,fontset=fandol]{ctexart}
\usepackage[a4paper,margin=1.78cm,headheight=15pt]{geometry}
\usepackage{amsmath,amssymb,mathtools,bm}
\usepackage{booktabs,tabularx,array}
\usepackage{enumitem}
\usepackage[most]{tcolorbox}
\usepackage{tikz}
\usetikzlibrary{arrows.meta,positioning}
\usepackage{xcolor}
\usepackage{fancyhdr}
\usepackage{hyperref}
\usepackage{bookmark}
\usepackage{microtype}

\newcolumntype{Y}{>{\raggedright\arraybackslash}X}
\newcolumntype{L}[1]{>{\raggedright\arraybackslash}p{#1}}
\setlength{\parindent}{2em}
\setlength{\parskip}{0.34em}
\setlength{\emergencystretch}{3em}
\renewcommand{\arraystretch}{1.2}
\setlength{\tabcolsep}{3pt}
\setlist[itemize]{itemsep=0.18em,topsep=0.35em,leftmargin=2em}
\setlist[enumerate]{itemsep=0.18em,topsep=0.35em,leftmargin=2em}

\definecolor{deep}{RGB}{42,58,73}
\definecolor{soft}{RGB}{245,247,249}
\definecolor{accent}{RGB}{104,76,55}
\definecolor{greenish}{RGB}{57,92,76}
\definecolor{warn}{RGB}{136,68,63}
\definecolor{purple}{RGB}{87,72,116}
\hypersetup{colorlinks=true,linkcolor=deep,urlcolor=accent,pdftitle={常微分方程：局部规律与整体轨迹}}
\pagestyle{fancy}
\fancyhf{}
\lhead{\small\color{deep}\textbf{数学一 · 高等数学 Topic 12}}
\rhead{\small\color{deep}结构识别 · 解空间 · 定解}
\cfoot{\small\thepage}
\renewcommand{\headrulewidth}{0.4pt}
\newtcolorbox{corebox}[1]{breakable,colback=soft,colframe=deep,title=\textbf{#1},boxrule=0.8pt,arc=2mm}
\newtcolorbox{methodbox}[1]{breakable,colback=white,colframe=greenish,title=\textbf{#1},boxrule=0.7pt,arc=1.5mm}
\newtcolorbox{warnbox}[1]{breakable,colback=white,colframe=warn,title=\textbf{#1},boxrule=0.7pt,arc=1.5mm}
\newtcolorbox{examplebox}[1]{breakable,colback=white,colframe=purple,title=\textbf{#1},boxrule=0.7pt,arc=1.5mm}
\newcommand{\kw}[1]{\textbf{\color{deep}#1}}

\begin{document}
\begin{center}
{\LARGE\textbf{常微分方程：局部规律与整体轨迹}}\\[0.4em]
{\large 从“变化率约束”到“解族与定解条件”}
\end{center}
\vspace{0.4em}
\begin{quote}
微分方程题的核心不是积分技巧本身，而是：已知局部变化率，如何识别隐藏结构、选择不改变解集的变换，并用初值或边界条件从解族中裁出唯一轨迹。
\end{quote}
\begin{center}\rule{0.5\linewidth}{0.5pt}\end{center}

\setcounter{section}{-1}
\section{本册定位}
本册承接 Topic03 的导数与局部变化率、Topic05 的积分和变限积分、Topic11 的级数型函数表示。主线覆盖数学一常见一阶方程、可降阶方程、常系数线性高阶方程、非齐次特解与共振、初值/边值和分段缝合。

\subsection{Own}
\begin{itemize}
  \item 方程阶数、线性/非线性、齐次/非齐次的结构语言；
  \item 可分离变量、一阶线性、齐次型、恰当、伯努利和整体代换；
  \item 不显含 $x$ 或 $y$ 的降阶；
  \item 高阶线性方程的解空间、基础解系、Wronskian 与“通解 = 齐次 + 特解”；
  \item 常系数特征根、非齐次待定系数、共振降阶与分段缝合；
  \item 初值/边值条件、定义区间、丢解和结果代回；
  \item 陌生方程中的自/因变量互换、独立变量重参数化、整体导数识别与算子因式降阶；
  \item 含变换自变量的函数—导数关系消元，包括对合与周期平移有限轨道，以及由解族反求微分方程；
  \item 周期强迫下一阶线性方程周期解的存在、唯一与退化边界。
\end{itemize}

\subsection{Uses / Stop Boundary}
\begin{itemize}
  \item Topic03 提供导数、链式法则和局部变化率；Topic05 提供原函数、变限积分与反常积分；
  \item Topic11 提供幂级数解的接口；B05 解释“特解 + kernel”的跨学科结构；
  \item 矩阵指数、相平面定性理论、非线性稳定性与数值 ODE 只作 Extension；
  \item 单一 ODE 问题族的题目触发信号、第一动作、检查点和停止条件进入同目录训练 Markdown；跨多个训练专题稳定复用并经证据验证的控制才进入《高等数学做题规则》。
\end{itemize}

\section{母模型：规律、变换与定解}
\[
\boxed{\begin{gathered}
\text{Equation Class}\to\text{Legal Transformation}\to\text{Integral / Kernel}\\
\to\text{General Solution}\to\text{Initial or Boundary Data}\to\text{Domain / Residual Check}
\end{gathered}}
\]
\begin{center}
\begin{tikzpicture}[
  node distance=0.52cm and 0.5cm,
  box/.style={draw=deep,rounded corners=1.5mm,minimum width=3.0cm,minimum height=0.92cm,align=center,fill=soft,font=\small},
  arr/.style={-{Latex[length=2mm]},thick,draw=deep}
]
\node[box] (a) {识别结构\\一阶/高阶};
\node[box,right=of a] (b) {合法变换\\代换/因子};
\node[box,right=of b] (c) {积分或算子\\恢复解族};
\node[box,below=of c] (d) {定解条件\\裁剪常数};
\node[box,left=of d] (e) {定义域与残差\\代回复核};
\draw[arr] (a)--(b); \draw[arr] (b)--(c); \draw[arr] (c)--(d); \draw[arr] (d)--(e);
\end{tikzpicture}
\end{center}

\section{第一关：一阶方程结构识别}
面对 $F(x,y,y')=0$，先问“哪个组合在重复出现”，而不是先把 $y'$ 当作可分离的符号。

\begin{tabularx}{\linewidth}{L{0.24\linewidth}L{0.31\linewidth}Y}
\toprule
类型 & 识别信号 & 第一动作 \\
\midrule
可分离 & $y'=g(x)h(y)$ & 先查 $h(y)=0$ 的常值解，再分离并积分 \\
一阶线性 & $y'+P(x)y=Q(x)$ & 积分因子 $\mu=e^{\int Pdx}$，化为 $(\mu y)'=\mu Q$ \\
齐次型 & $y'=F(y/x)$ 或同次 $Mdx+Ndy=0$ & 令 $y=vx$，$y'=v+xv'$ \\
恰当 & $Mdx+Ndy=0$ 且 $M_y=N_x$ & 求势函数 $\Phi_x=M,\Phi_y=N$ \\
伯努利 & $y'+Py=Qy^\alpha$ & 令 $z=y^{1-\alpha}$，先保留 $y=0$ 等退化支 \\
整体代换 & 反复出现 $ax+by+c$ 或组合量 & 令 $u=ax+by+c$，同步替换导数与自变量 \\
\bottomrule
\end{tabularx}

\begin{corebox}{恰当方程与路径无关其实共享同一个势函数}
若
\[
M(x,y)\,dx+N(x,y)\,dy=0
\]
是恰当方程，存在势函数 $\Phi$ 使
\[
d\Phi=M\,dx+N\,dy.
\]
于是微分方程的解曲线就是势函数的等值线
\[
\boxed{\Phi(x,y)=C}.
\]
同一个结构放到 Topic09 看，就是向量场 $(M,N)=\nabla\Phi$ 的第二型曲线积分只依赖端点；在单连通且 $C^1$ 的常规考试区域内，$M_y=N_x$ 正是把局部相容升级为全局势函数的常用资格。

因此“恰当方程”和“路径无关”不是两套公式，而是同一个势函数机制在两个问题里的不同输出：前者问\textbf{哪些曲线保持势值不变}，后者问\textbf{沿任意路径势值改变了多少}。若定义域有孔洞，$M_y=N_x$ 仍只给局部相容，不能跳过 Topic09 的拓扑资格审计。
\end{corebox}

\begin{methodbox}{一阶线性方程的积分因子}
\[
y'+P(x)y=Q(x),\qquad \mu(x)=e^{\int P(x)\,dx},
\]
则
\[
(\mu y)'=\mu Q,\qquad
y=\frac{1}{\mu(x)}\left(C+\int \mu(x)Q(x)\,dx\right).
\]
变限形式
\[
y(x)=e^{-\int_{x_0}^{x}P(t)\,dt}\left[y_0+\int_{x_0}^{x}e^{\int_{x_0}^{s}P(t)\,dt}Q(s)\,ds\right]
\]
天然吸收初值，适合做存在性与稳定性检查。
\end{methodbox}

\begin{warnbox}{变形不改变解集是硬条件}
除以 $y,h(y),x$ 或最高阶系数之前，先检查它们为零的分支；平方、开方和取对数后要回代。解必须连同定义区间报告，初值点不能落在系数奇点上。
\end{warnbox}

\section{第二关：降阶与一阶非线性}
若二阶方程不显含 $y$，令 $p=y'$，则 $y''=dp/dx$；若不显含 $x$，令 $p(y)=y'$，则 $y''=p\,dp/dy$。降阶后回到一阶结构，积分两次并补足常数。

对齐次型 $y'=F(y/x)$，令 $v=y/x$ 后得到 $y'=v+xv'$，通常可分离；直线解 $y=Cx$ 需要在代换前后单独检查。伯努利方程的代换把非线性幂变成线性，但 $y=0$ 等被除掉的支不能遗漏。

\section{第三关：陌生方程先改表示，不要急着给它贴类型标签}

标准类型识别失败时，优先问：\textbf{当前选择的自变量、因变量或状态量是否让结构变复杂了？} 只要变换在当前区间合法，换表示并不改变轨迹本体。

\subsection{交换自变量与因变量}

若解曲线局部满足 $y'(x)\ne0$，可把它改写为 $x=x(y)$。设
\[
p(y)=\frac{dx}{dy},
\]
则
\[
\frac{dy}{dx}=\frac1p,
\qquad
\frac{d^2y}{dx^2}=-\frac{p'}{p^3}
=-\frac{d^2x/dy^2}{(dx/dy)^3}.
\]
因此某些对 $y'$ 呈高次、但对 $dx/dy$ 呈线性的方程，会在交换变量后降复杂度。资格边界是局部可逆：若 $y'=0$，不能无条件使用倒数公式，必须另查该分支。

\subsection{重参数化独立变量：先搬运导数算子，再看方程是否降复杂度}

有些方程并不需要交换自变量与因变量，只需要把同一条解曲线改用新的参数描述。设
\[
x=\phi(t),\qquad y(x)=Y(t),\qquad \phi'(t)\ne0,
\]
则
\[
\frac{dy}{dx}=\frac{Y_t}{\phi_t},
\qquad
\frac{d^2y}{dx^2}
=\frac1{\phi_t}\frac d{dt}\left(\frac{Y_t}{\phi_t}\right)
=\frac{Y_{tt}\phi_t-Y_t\phi_{tt}}{\phi_t^3}.
\]
若反过来直接设新变量 $t=\psi(x)$，则
\[
y'=Y_t\psi',
\qquad
y''=Y_{tt}(\psi')^2+Y_t\psi''.
\]
因此 $t=e^x$ 时有 $y'=tY_t,\ y''=t^2Y_{tt}+tY_t$；$x=\sin t$ 时有
\[
y'=\frac{Y_t}{\cos t},
\qquad
y''=\frac{Y_{tt}\cos t+Y_t\sin t}{\cos^3 t}.
\]
这类换元的资格不是“题目给了一个漂亮代换”就自动成立，而是新参数在当前区间必须局部一一；例如 $x=\sin t$ 需要避开 $\cos t=0$ 的退化点，并在单调分支上工作。

\subsection{识别“整个导数”作为新状态}

若反复出现
\[
xy'+y=(xy)',
\qquad
\frac{xy'-y}{x^2}=\left(\frac yx\right)',
\qquad
\frac{yy''-(y')^2}{y^2}=\left(\frac{y'}y\right)',
\]
则优先把括号中的整体设为新状态。其本质不是凑公式，而是寻找一个使方程只依赖较少状态变量的表示。

\subsection{算子因式分解提供系统降阶}

对常系数算子 $D=d/dx$，若
\[
L(D)=L_1(D)L_2(D),
\]
可先设 $u=L_2(D)y$，求低阶方程 $L_1(D)u=f$，再解 $L_2(D)y=u$。特别地，
\[
(D-a)(e^{ax}v)=e^{ax}v',
\]
所以
\[
(D-a)^m y=e^{ax}v^{(m)}\qquad (y=e^{ax}v).
\]
这解释了为何像 $(D-1)^2y=e^x/x$ 这类右端不适合待定系数的方程，仍可通过 $y=e^xv$ 直接降成 $v''=1/x$。

\section{第四关：高阶线性方程的解空间}
线性方程写成
\[
L[y]=a_n(x)y^{(n)}+\cdots+a_1(x)y'+a_0(x)y=f(x).
\]
在系数连续、最高阶系数不为零的区间上，齐次解集是 $n$ 维向量空间；给定一组合法的 $n$ 个初值，局部解唯一。非齐次解集是一个仿射平移：
\[
\boxed{y=y_h+y_p},
\]
其中 $y_h$ 解 $L[y]=0$，$y_p$ 是任一特解。

二阶齐次方程的基础解系 $y_1,y_2$ 可用 Wronskian
\[
W(y_1,y_2)=\begin{vmatrix}y_1&y_2\\y_1'&y_2'\end{vmatrix}
\]
检查线性无关。不要把“二阶”机械理解为任意情况下都恰有两个自由常数；退化点和奇异区间先过资格门。

\section{第五关：常系数、特征根与共振}
对 $a_ny^{(n)}+\cdots+a_0y=0$，试探 $y=e^{rx}$ 得特征方程
\[
a_nr^n+\cdots+a_1r+a_0=0.
\]
\begin{tabularx}{\linewidth}{L{0.28\linewidth}Y}
\toprule
根型 & 齐次解块 \\
\midrule
不同实根 $r_1,r_2$ & $C_1e^{r_1x}+C_2e^{r_2x}$ \\
二重实根 $r$ & $(C_1+C_2x)e^{rx}$ \\
共轭根 $\alpha\pm\beta i$ & $e^{\alpha x}(C_1\cos\beta x+C_2\sin\beta x)$ \\
\bottomrule
\end{tabularx}

非齐次右端 $g(x)$ 先按指数、多项式、正弦余弦及其乘积分解，再猜特解。若试探项已出现在齐次解中，乘以 $x^s$，其中 $s$ 是对应特征根的重数；这就是共振修正。

\begin{examplebox}{贯穿母例：共振方程从结构识别到定解}
求解
\[
y''-2y'+y=e^x,\qquad y(0)=0,\quad y'(0)=0.
\]
算子为 $L[D]=(D-1)^2$，齐次解是 $(C_1+C_2x)e^x$。右端 $e^x$ 对应二重根 $r=1$，故普通试探 $Ae^x$ 或 $Axe^x$ 都落在齐次核中；乘足 $x^2$，取 $y_p=Ax^2e^x$，代回得 $2Ae^x=e^x$，所以 $A=1/2$。通解
\[
y=e^x\left(C_1+C_2x+\frac{x^2}{2}\right)
\]
经两个初值裁剪为 $C_1=C_2=0$。最后直接代回原方程得到残差 $e^x$，并核对两个初值。

这条轨迹完整经过“结构识别 $\to$ 合法试探 $\to$ 齐次空间加特解 $\to$ 定解裁剪 $\to$ 残差复核”。最典型的第一次偏离是共振时没有乘够 $x$ 的次数，导致试探项仍被齐次算子消掉。
\end{examplebox}

\section{第六关：Euler--Cauchy 方程把尺度不变性翻译成常系数}

形如
\[
x^2y''+axy'+by=g(x)\qquad(x>0\text{ 或 }x<0)
\]
的 Euler--Cauchy 方程看似变系数，但系数恰好按导数阶数补偿缩放。它的母结构是\textbf{乘法尺度 $x\mapsto cx$ 下保持同型}，因此最自然的坐标不是 $x$，而是对数时间
\[
t=\ln\lvert x\rvert.
\]
令 $Y(t)=y(e^t)$（在 $x>0$ 分支），则
\[
xy'=Y_t,
\qquad
x^2y''=Y_{tt}-Y_t,
\]
从而原方程被翻译成常系数方程
\[
Y_{tt}+(a-1)Y_t+bY=g(e^t).
\]

齐次情形也可直接试探 $y=x^m$：
\[
x^2y''+axy'+by
=x^m\bigl[m(m-1)+am+b\bigr],
\]
于是得到关于幂指数 $m$ 的特征方程。$t=\ln\lvert x\rvert$ 解释了为什么幂函数对应常系数世界中的指数函数 $e^{mt}$；直接试 $x^m$ 是这一坐标变换的压缩版。

\begin{tcolorbox}[title=Euler--Cauchy 的资格边界]
对数换元必须固定在不穿过 $x=0$ 的连通分支上；$x=0$ 是方程的奇异点，不能把 $x>0$ 与 $x<0$ 的解族用同一组常数无条件跨接。
\end{tcolorbox}

\section{第七关：函数变换关系与积分关系怎样压成局部 ODE}

\subsection{变换自变量的函数—导数关系：求导、回写、消元}

若
\[
f'(x)=f(\phi(x)),
\]
右侧采样的是另一个自变量位置。直接把它当普通自治方程通常无效；应对关系求导：
\[
f''(x)=f'(\phi(x))\phi'(x).
\]
若变换满足 $\phi(\phi(x))=x$（对合），则原关系在 $\phi(x)$ 处给出
\[
f'(\phi(x))=f(x),
\]
从而闭合为
\[
\boxed{f''(x)=\phi'(x)f(x)}.
\]
例如 $\phi(x)=1-x$ 得 $f''+f=0$；$\phi(x)=2/x$（$x>0$）得
\[
x^2f''+2f=0.
\]

对合只是“采样位置经过有限步后回到原点”的最短情形。若还给周期性，平移也可能形成有限轨道。例如 $f$ 为 $2\pi$ 周期且
\[
f(x)+f'(x+\pi)=\sin x,
\]
平移 $x\mapsto x+\pi$ 两次就回到同一采样位置。把原式平移 $\pi$ 得
\[
f(x+\pi)+f'(x)=-\sin x.
\]
对这条式子求导，再用原式回写 $f'(x+\pi)=\sin x-f(x)$，得到
\[
\boxed{f''(x)-f(x)=-\sin x-\cos x}.
\]
一般地，若平移量 $a$ 与周期 $T$ 满足 $na=mT$，则 $x,x+a,\ldots,x+(n-1)a$ 只形成有限个采样状态，可以尝试“平移/求导—回写—消元”；若 $a/T$ 无理，采样轨道通常不会有限闭合，不能机械套这条路线。

消元后得到的 ODE 是原函数关系的必要后果；求出 ODE 解后仍必须代回原函数关系，因为微分与消元可能扩大解集。

\subsection{从积分/函数关系反推 ODE——先微分消掉累积，再补回丢失条件}

有些题目不给 ODE，而给函数关系，例如
\[
y(x)=c+\int_{x_0}^{x}K(x,t,y(t))\,dt
\]
或含若干次变上限积分、未知积分常数。此时“对等式求导”不是纯计算技巧，而是在做\textbf{状态恢复}：每求一次导数，就消掉一层累积记忆，直到暴露出局部微分规律。

最简单的线性例子：
\[
y(x)=A+B(x-x_0)+\int_{x_0}^{x}(x-t)f(t)\,dt.
\]
由 H-B03 的重复累积接口，连续求导得到
\[
y'(x)=B+\int_{x_0}^{x}f(t)\,dt,
\qquad
y''(x)=f(x).
\]
原积分关系同时自动编码初值
\[
y(x_0)=A,\qquad y'(x_0)=B.
\]
因此从积分方程转 ODE 后，必须把这些被微分“抹掉”的低阶信息重新作为初值带回去；只求出 ODE 通解而忘记恢复原关系，会引入额外伪解。

同理，若函数方程中出现一个未知整体量
\[
C=\int_a^b y(t)\,dt,
\]
可先把 $C$ 当作标量参数解局部 ODE，再把所得 $y$ 代回积分关系闭合 $C$。这和 Topic08 的“未知整体积分先设常数”是同一种\textbf{低维闭环}机制。

\section{第八关：由解族反求方程——自由参数是需要被消去的隐藏状态}

若给出含任意常数的解族，反求方程的核心是通过足够次数求导消去这些常数。一个独立任意常数通常对应一阶关系，$n$ 个真正独立的常数通常需要到 $n$ 阶；但必须警惕参数冗余或退化，不能只数符号个数。

例如
\[
(x+C)^2+y^2=1
\]
对 $x$ 求导得
\[
2(x+C)+2yy'=0.
\]
由原式与导数式消去 $C$，得到不含任意常数的微分关系。反求完成后应检查原解族逐个代回确实成立。

对常系数线性方程，若已知通解中的齐次块，可以从 $e^{rx}$、$xe^{rx}$、$e^{\alpha x}\cos\beta x$ 等模式恢复特征根及重数；若还给出一个特解 $y_p$，则用恢复出的齐次算子 $L(D)$ 计算 $L(D)y_p$，得到对应非齐次右端。这里“解的结构”比逐项猜系数更稳定。

\section{第九关：周期解、定解、分段缝合与稳定性接口}

对一阶线性方程
\[
y'+ay=f(x),
\]
若 $f$ 连续且以 $T$ 为周期，变限公式给
\[
y(x)=e^{-ax}\left(C+\int_0^x e^{at}f(t)\,dt\right).
\]
周期条件 $y(T)=y(0)$ 给
\[
(1-e^{-aT})C=e^{-aT}\int_0^T e^{at}f(t)\,dt.
\]
为什么单点端值条件 $y(T)=y(0)$ 足以保证函数全域具有周期性？令平移解 $z(x)=y(x+T)$，由 $f(x+T)=f(x)$ 知 $z'(x)+az(x)=f(x+T)=f(x)$，即 $z$ 与 $y$ 满足完全相同的 ODE；又初值 $z(0)=y(T)=y(0)$，由一阶线性方程初值解的唯一性（Picard）立即得到 $z(x)\equiv y(x)$，即 $y(x+T)\equiv y(x)$ 全局成立。

因此对实常数 $a\ne0$，存在且仅存在一个 $T$ 周期解；若 $a=0$，方程退化为 $y'=f(x)$，周期解存在当且仅当
\[
\int_0^T f(t)\,dt=0,
\]
且一旦存在便可任加常数，所以不唯一。这揭示“周期强迫必有唯一周期响应”并非无条件事实，关键是周期边界条件对齐次自由度能否唯一裁剪。

随后再处理一般初值、边值与分段缝合：
初值条件用于确定通解常数；边界条件要把所有候选解同时代入，可能无解或多解。分段右端项时，在切换点分别求左右解，再按题设连续性/可导性缝合；不能把一个区间的常数直接沿用到另一区间。

对一维自治方程
\[
y'=\varphi(y),
\]
平衡点 $y_*$ 满足 $\varphi(y_*)=0$。真正决定局部稳定性的不是“解出平衡点”本身，而是附近箭头方向：
\[
\varphi(y)>0\Rightarrow y\text{ 向上运动},
\qquad
\varphi(y)<0\Rightarrow y\text{ 向下运动}.
\]
若 $y_*$ 左侧 $\varphi>0$、右侧 $\varphi<0$，两边轨道都指向 $y_*$，故局部稳定；符号相反则不稳定。若 $\varphi$ 可导且 $\varphi'(y_*)\ne0$，这压缩为
\[
\varphi'(y_*)<0\Rightarrow\text{局部吸引},
\qquad
\varphi'(y_*)>0\Rightarrow\text{局部排斥}.
\]
当 $\varphi'(y_*)=0$ 时线性化失效，必须回到原函数符号或更高阶项，不能把“导数为零”当作稳定。

对二阶常系数齐次方程，特征根实部控制指数包络：实部为负时自由响应衰减，为正时增长；\textbf{简单纯虚根}给出有界的持续振荡。若零实部根有重数，$x^k e^{i\omega x}$ 的多项式因子会造成无界增长，因此不能只看到“实部为零”就宣布稳定振荡。周期强迫项在共振时会出现额外的 $x$ 因子；要存在与强迫同周期的有界稳态解，必须避免相应频率落入齐次核导致的共振增长。完整线性系统稳定性、Routh--Hurwitz 等只作 Extension，不进入数学一主干。

\begin{corebox}{求解结束的五项复核}
\begin{enumerate}
  \item 方程阶数、线性/非线性和齐次含义是否识别正确；
  \item 变换是否丢掉了零解、直线解或奇异支；
  \item 通解自由常数个数是否与解空间维数匹配；
  \item 初值/边值是否全部代入且定义区间合法；
  \item 将结果代回原方程，检查残差、连续性和共振次数。
\end{enumerate}
\end{corebox}

\section{诊断流程与证据回链}
\begin{center}
\begin{tikzpicture}[
  node distance=0.42cm and 0.55cm,
  box/.style={draw=deep,rounded corners=1mm,minimum width=3.3cm,minimum height=0.8cm,align=center,fill=soft,font=\small},
  arr/.style={-{Latex[length=1.8mm]},thick,draw=deep}
]
\node[box] (a) {识别阶数与类型};
\node[box,below=of a] (b) {合法代换/积分因子};
\node[box,below=of b] (c) {积分或特征根};
\node[box,right=of c] (d) {齐次解 + 特解};
\node[box,right=of d] (e) {初值/边值裁剪};
\node[box,above=of e] (f) {定义域/残差复核};
\draw[arr] (a)--(b); \draw[arr] (b)--(c); \draw[arr] (c)--(d); \draw[arr] (d)--(e); \draw[arr] (e)--(f);
\end{tikzpicture}
\end{center}

\begin{tabularx}{\linewidth}{L{0.28\linewidth}Y}
\toprule
来源段落 & 本册保留的机制与边界 \\
\midrule
I-13 §1--§4 & 阶数、线性/非线性、解族、定解条件、丢解与定义域 \\
I-13 §5--§10 & 一阶结构识别、分离、积分因子、齐次型、恰当、伯努利、整体代换 \\
I-13 §13--§24 & 降阶、高阶线性解空间、Wronskian、非齐次仿射结构 \\
I-13 §25--§37 & 常系数特征根、待定系数、共振与算子接口 \\
I-13 §38--§55 & 分段缝合、Euler--Cauchy 的对数尺度变换、积分/函数关系反推 ODE 与初值恢复、自治相线和常系数稳定性（主干 + Extension 边界） \\
《高数下册进阶》专题3 & 自/因变量互换、整体导数、算子因式降阶、变换自变量消元、由解族反求方程与周期解资格；跨章节构造只保留“翻译成 ODE”接口，原章节机制仍归各 Topic \\
\bottomrule
\end{tabularx}

\begin{corebox}{一句话压缩}
\textbf{先辨认变化率方程的结构，再用不丢解的变换恢复解族；高阶线性统一为“齐次空间 + 特解”，最后由初值、边值和定义区间裁出可用轨迹。}
\end{corebox}

\section{陌生题验证协议}
\begin{enumerate}
  \item 写出阶数、线性/非线性、齐次/齐次型的判定依据；
  \item 变换前列出可能为零的因子，变换后回代特殊支；
  \item 一阶题先走结构表，高阶线性先求特征根和齐次解；
  \item 非齐次题按右端类型猜特解，遇到共振乘足够次数的 $x$；
  \item 初值/边值全部代入，分段题在接口点执行连续性/可导性缝合；
  \item 最后做定义域、残差、常数个数和增长/衰减包络复核。
\end{enumerate}

\end{document}
